\section{5. Ground Segment Life
Profile}\label{ground-segment-life-profile}

The life profile covers design, integration, verification,
commissioning, operations and retirement of terrestrial infrastructure,
coordinated with the SSEG lifecycle.

\subsection{5.1 Development and
Preparation}\label{development-and-preparation}

Requirements, interfaces and the GSEG configuration are defined,
followed by preparation of infrastructure, software, procedures and
operational data. The seven subsystems cover communications,
constellation control, Flight Dynamics, processing, warning, archiving
and common infrastructure. Development is coordinated with SSEG for
exchanged products and mission functions.

\subsection{5.2 Integration, Verification and Operational
Validation}\label{integration-verification-and-operational-validation}

Installation and configuration of equipment and integration of the seven
subsystems are covered. Technical verification checks functions,
performance and interfaces against requirements and the verification
programme. Operational validation concerns personnel, procedures, data
and coordination with SSEG. Integration tests, representative
simulations, HIL and complete warning-chain exercises
(\texttt{R-\allowbreak{}038}--\texttt{R-\allowbreak{}040}) are employed as applicable.

\subsection{5.3 Launch Support and Constellation
Commissioning}\label{launch-support-and-constellation-commissioning}

GSEG achieves readiness before spacecraft operations needing ground
support. During launch and LEOP it provides communications, telemetry,
commands and supervision in coordination with SSEG. During GEO and LEO
commissioning it confirms communication interfaces, availability of data
products and warning-function readiness.

\newpage
\subsection{5.4 Operation and
Maintenance}\label{operation-and-maintenance}

GSEG provides continuous constellation supervision and control, data
reception and analysis, and warning dissemination. GIS monitors and
maintains infrastructure, while MCC supervises the mission and plans
orbital operations. Maintenance covers diagnosis, configuration
management, repair and controlled updates; changes are verified before
deployment. Security, redundancy and availability are considered
(\texttt{R-\allowbreak{}015}, \texttt{R-\allowbreak{}017}, \texttt{R-\allowbreak{}030}--\texttt{R-\allowbreak{}037}).

\subsection{5.5 End of Mission and GSEG
Decommissioning}\label{end-of-mission-and-gseg-decommissioning}

GSEG supports spacecraft commands and monitoring during mission
close-out; spacecraft disposal remains an SSEG responsibility
(\texttt{R-\allowbreak{}018}). This includes terminating ground services, retaining
documentation and data, and retiring equipment and software. Archived
data remain subject to protection, access and retention requirements.

\begingroup\GSEGTableStyle{3}
\begin{longtable}[]{|P{0.19}|P{0.49}|P{0.32}|}
\hline\rowcolor{TableHeaderBlue}
\begin{minipage}[b]{\linewidth}\raggedright
\color{white}\bfseries\sffamily
Phase
\end{minipage} & \begin{minipage}[b]{\linewidth}\raggedright
\color{white}\bfseries\sffamily
Principal activities
\end{minipage} & \begin{minipage}[b]{\linewidth}\raggedright
\color{white}\bfseries\sffamily
End state
\end{minipage} \\
\hline
\endhead
\hline
\endlastfoot
1. Development & Requirements, design, software, infrastructure and
operational-data preparation. & Ready for integration. \\ \hline
2. Integration and validation & Installation, integration, verification
and operational exercises. & Ground readiness confirmed. \\ \hline
3. Launch and commissioning & Launch/LEOP support, communication checks
and spacecraft commissioning. & Ready for nominal operations. \\ \hline
4. Routine operations & Mission operations, maintenance, updates and
fault handling. & Operational capability maintained. \\ \hline
5. Decommissioning & End-of-mission support, records retention and
retirement of services. & GSEG operations ended. \\*
\end{longtable}
\endgroup
